
\SGASection{7}{Cartan subgroups of a smooth group}
\SGARefTarget{sga3:xv:section:7}

\SGAStatement{\SGAInlineRefTarget{sga3:xv:proposition:7-1}Proposition 7.1}{Let \(S\) be a prescheme and let \(G\)
be an \(S\)-group prescheme of finite presentation, smooth over \(S\),
with connected fibres.
\begin{enumerate}
\renewcommand{\labelenumi}{\textup{\roman{enumi})}}
\item\SGAInlineRefTarget{sga3:xv:proposition:7-1:item:i} Let \(\mathcal{CT}\) be the \(S\)-functor defined by
\[
\mathcal{CT}(S')=
\left\{
\begin{array}{c}
\text{subgroup preschemes \(H\) of \(G_{S'}\), smooth over \(S'\),}\\
\text{such that every \(H_{s'}\) is the centralizer in \(G_{s'}\)}\\
\text{of a sub-torus of \(G_{s'}\)}
\end{array}
\right\}.
\]
Then \(\mathcal{CT}\) is represented by an \(S\)-prescheme of finite
presentation, smooth and quasi-projective over \(S\).
\item If \(S\) is artinian and \(H\in\mathcal{CT}(S)\), then \(H\)
is the centralizer in \(G\) of a sub-torus of \(G\). If \(S\) is the
spectrum of a henselian local ring and \(H\in\mathcal{CT}(S)\), then
\(H\) is the centralizer in \(G\) of a subgroup of multiplicative
type that is étale over \(S\).
\item\SGAInlineRefTarget{sga3:xv:proposition:7-1:item:iii} Let \(L\) be an \(S\)-group prescheme smooth and of finite
presentation, let \(i:L\to G\) be an \(S\)-group monomorphism, and
let \(H\in\mathcal{CT}(S)\). Then
\(\operatorname{Transp}_G(H,L)\) and
\(\operatorname{Transpstr}_G(H,L)\)
(Exposé \SGARefLink{sga3:viii:examples:6:5:item:e:01}{VIII~6.5(e)}) are represented by closed subpreschemes of \(G\),
smooth over \(S\).
\item If \(H\in\mathcal{CT}(S)\), then \(H\) is closed in \(G\);
\(N=\operatorname{Norm}_G(H)\) is represented by a closed subgroup
prescheme of \(G\), smooth over \(S\); \(N/H\) is represented by an
\(S\)-group prescheme, separated, étale, and of finite type over
\(S\); and \(G/N\) is represented by an \(S\)-prescheme smooth and
quasi-projective over \(S\).
\item\SGAInlineRefTarget{sga3:xv:proposition:7-1:item:v} Let \(G'\) be an \(S\)-group prescheme of finite presentation
and let \(u:G\to G'\) be faithfully flat. Then \(G'\) satisfies the
same hypotheses as \(G\) (Exposé VI\textsubscript{B}~9). If
\(H\in\mathcal{CT}_G(S)\), its image is represented by a subgroup
prescheme \(H'\in\mathcal{CT}_{G'}(S)\), and \(H\to H'\) is faithfully
flat. If \(H=\operatorname{Centr}_G(T)\) for a torus \(T\), then
\[
 H'=\operatorname{Centr}_{G'}(T'),\qquad T'=u(T).
\]
\item\SGAInlineRefTarget{sga3:xv:proposition:7-1:item:vi} Under the hypotheses of (v), the \(S\)-morphism
\[
 \widetilde u:\mathcal{CT}_G\longrightarrow\mathcal{CT}_{G'},
 \qquad H\longmapsto H'=u(H),
\]
is quasi-compact and faithfully flat. If, moreover, \(\ker u\) is
central, then \(\widetilde u\) is an isomorphism, with inverse
\(H'\mapsto u^{-1}(H')\).
\end{enumerate}}
\SGARefTarget{sga3:xv:proposition:7:1}

\medskip\noindent\textit{Proof of (ii).}
For the first assertion we may suppose that \(S\) is local artinian,
with closed point \(s\). Let \(H\in\mathcal{CT}(S)\), and let \(T_s\)
be the maximal central torus of \(H_s\), which is defined over
\(\kappa(s)\) (cf. \SGARefLink{sga3:xv:lemma:3:4}{Lemma 3.4}). Since
\[
  H_s=\operatorname{Centr}_{G_s}(T_s),
\]
and \(H\) is smooth, \(T_s\) lifts uniquely to a sub-torus \(T\) of
\(H\), central in \(H\) (Exposé \SGARefLink{sga3:ix:theorem:3:6:bis}{IX~3.6 bis} and \SGARefLink{sga3:xv:remark:5:6}{5.6}). Then
\(H'=\operatorname{Centr}_G(T)\) contains \(H\) and has the same
fibre as \(H\). Because \(H\) is flat, \(H=H'\) (Exposé
VI\textsubscript{B}~\SGARefLink{sga3:xv:lemma:2:5}{2.5}).

Suppose now that \(S\) is the spectrum of a henselian local ring,
which the usual reductions allow us to take noetherian. Let \(s\) be
its closed point, let \(T_s\) be the maximal central torus of \(H_s\),
let \(q\) be invertible on \(S\), let \(\ell\) be a power of \(q\),
and let \(T\) be an \(S\)-torus whose closed fibre is isomorphic to
\(T_s\) (Exposé X~4.6). Let \({}_\ell H\) denote the ``kernel'' of
the \(\ell\)-th power morphism in \(H\), and let \(U_\ell\) be the
largest open subset of \({}_\ell H\) étale over \(S\). \SGARefLink{sga3:xv:lemma:1:3}{Lemma 1.3} and
flatness of \(H\) show that \(U_\ell\) contains \({}_\ell T_s\).
Hensel's lemma gives a unique \(S\)-morphism
\[
  u_\ell:{}_\ell T\longrightarrow U_\ell
\]
inducing the canonical immersion
\({}_\ell T_s\to(U_\ell)_s\) on the closed fibre. Uniqueness shows
that \(u_\ell\) is a central homomorphism (Exposé \SGARefLink{sga3:ix:theorem:5:6:item:a}{IX~5.6(a)}).
Proceeding as in \SGARefLink{sga3:xv:lemma:6:6}{6.6} and \SGARefLink{sga3:xv:lemma:6:7}{6.7}, one proves that \(u_\ell\) is an
immersion and, if \(M_\ell=u_\ell({}_\ell T)\), that
\[
  \operatorname{Centr}_G(M_\ell)=H
\]
for all sufficiently large \(\ell\). This is where noetherianity of
\(S\) is used.

\medskip\noindent\textit{Proof of (i).}
If \(H\in\mathcal{CT}(S)\), then \(H\) has connected fibres (Exposé
\SGARefLink{sga3:xii:theorem:6-6:item:b}{XII~6.6(b)}) and equals its connected normalizer (\SGARefLink{sga3:xv:remark:6:5}{6.5}). Hence the
present functor \(\mathcal{CT}\) is the functor of \SGARefLink{sga3:xv:paragraph:6-4-0}{6.4.0}. \SGARefLink{sga3:xv:theorem:6:4}{Theorem 6.4}
therefore represents it by an \(S\)-prescheme of finite presentation,
quasi-projective over \(S\). It remains to prove smoothness.

EGA IV~8 reduces us to affine noetherian \(S\). By Exposé XI~1.5 it
then suffices to prove the following infinitesimal lifting property:
if \(S\) is local artinian, \(S_0\) is defined by a nilpotent ideal,
and \(H_0\in\mathcal{CT}(S_0)\), then \(H_0\) lifts to a smooth
subgroup prescheme \(H\) of \(G\). By part (ii),
\[
 H_0=\operatorname{Centr}_{G_0}(T_0)
\]
for a sub-torus \(T_0\) of \(G_0\). Since \(G\) is smooth, \(T_0\)
lifts to a sub-torus \(T\) of \(G\) (Exposé \SGARefLink{sga3:ix:theorem:3:6:bis}{IX~3.6 bis}), and
\(H=\operatorname{Centr}_G(T)\) works; it is smooth by Exposé \SGARefLink{sga3:xi:corollary:2:4}{XI~2.4}
and \SGARefLink{sga3:xv:lemma:2:5}{Lemma 2.5}.

\medskip\noindent\textit{Proof of (iii).}
Because \(H\) is smooth with connected fibres, Exposé \SGARefLink{sga3:xi:corollary:6:11}{XI~6.11}
represents \(\operatorname{Transp}_G(H,L)\) by a closed subprescheme
of \(G\), of finite presentation over \(S\). Smoothness reduces to
the following lifting problem. Let \(S\) be local artinian, \(S_0\)
a closed subprescheme, and \(g_0\in G(S_0)\) such that
\[
  \operatorname{int}(g_0)H_0\subseteq L_0.
\]
We must lift \(g_0\) to \(g\in G(S)\) with
\(\operatorname{int}(g)H\subseteq L\).

Smoothness of \(G\) gives \(g_1\in G(S)\) lifting \(g_0\). Put
\(H'=\operatorname{int}(g_1)H\). Then
\(H'\in\mathcal{CT}(S)\) and \(H'_0\subseteq L_0\). By (ii),
\(H'=\operatorname{Centr}_G(T')\) for a torus \(T'\subseteq G\).
Since \(L\) is smooth, \(T'_0\subseteq L_0\) lifts to a torus
\(T''\subseteq L\) (Exposé \SGARefLink{sga3:ix:theorem:3:6:bis}{IX~3.6 bis}). The group
\(\operatorname{Centr}_L(T'')\) lies in
\(\operatorname{Centr}_G(T'')\), has the same closed fibre as the
latter, namely \(H'_0\), and is smooth. Hence
\[
 \operatorname{Centr}_L(T'')
 =\operatorname{Centr}_G(T'').
\]
The sub-tori \(T'\) and \(T''\) of \(G\) lift the same \(T'_0\), so
they are conjugate by an element \(h\in G(S)\) reducing to the
identity (Exposé \SGARefLink{sga3:ix:corollary:3:3:bis}{IX~3.3 bis}). Their centralizers are therefore
conjugate. Then \(g=hg_1\) lifts \(g_0\) and satisfies
\(\operatorname{int}(g)H\subseteq L\).

For \(g\in\operatorname{Transp}_G(H,L)(S)\), the equality
\(\operatorname{int}(g)H=L\) holds if and only if
\(\dim H_s=\dim L_s\) for every \(s\in S\). If \(U\) is the
open-and-closed subprescheme over which these dimensions agree, then
\[
 \operatorname{Transpstr}_G(H,L)
 =U\times_S\operatorname{Transp}_G(H,L).
\]

\medskip\noindent\textit{Proof of (iv).}
To prove \(H\) closed in \(G\), reduce first to affine noetherian
\(S\), then, by EGA IV~8, to the spectrum of a complete local ring.
Part (ii) writes \(H\) as the centralizer in \(G\) of a subgroup of
multiplicative type, so \(H\) is closed because \(G\) is separated
(Exposé VI\textsubscript{B}~\SGARefLink{sga3:xv:theorem:5-2}{5.2}).

By (iii),
\[
 N=\operatorname{Norm}_G(H)
  =\operatorname{Transpstr}_G(H,H)
\]
is represented by a subgroup prescheme smooth over \(S\) and closed
in \(G\). The morphism
\[
  G\longrightarrow\mathcal{CT},
  \qquad g\longmapsto\operatorname{int}(g)H
\]
is smooth by (iii), so its image is an open subset \(U\) of
\(\mathcal{CT}\). As in Exposé \SGARefLink{sga3:xi:corollary:5:3}{XI~5.3}, \(G/N\) is represented by
\(U\); it is therefore quasi-projective.

We now treat \(N/H\). By EGA IV~8, assume \(S\) affine noetherian and
then \(S=\operatorname{Spec}A\), with \(A\) local of finite
dimension\footnote{Krull dimension.}. We use induction on
\(\dim S\). When \(\dim S=0\), the assertion is Exposé
VI\textsubscript{A}, §4. If \(N/H\) is representable, it is separated
because \(H\) is closed in \(N\), of finite type and étale because
\(N\) is smooth of finite type and \(H\) is open. It is therefore
quasi-affine (SGA 1, VIII~6.2). Effective descent for quasi-affine
schemes (\emph{loc. cit.}, 7.9) permits replacement of \(A\) by its
completion.

Thus assume \(S\) complete noetherian local, with closed point \(s\),
and put \(U=S\setminus\{s\}\). By induction,
\((N|_U)/(H|_U)\) is represented by a \(U\)-group \(K\). Let \(T_s\)
be the maximal central torus of \(H_s\), choose \(q\) invertible on
\(S\), let \(\ell\) be a power of \(q\), and let \(M_\ell\) be the
unique central subgroup of multiplicative type of \(H\) lifting
\({}_\ell T_s\), as in (ii). Choose \(\ell\) so large that
\[
  \operatorname{Centr}_G(M_\ell)=H,
\]
and put \(N'=\operatorname{Norm}_G(M_\ell)\).

The torus \(T_s\), and hence \((M_\ell)_s\), is characteristic in
\(H_s\). Therefore \(N_s\) normalizes \((M_\ell)_s\), and
\(N_s=N'_s\). The proof of Exposé \SGARefLink{sga3:xi:proposition:5:9}{XI~5.9} represents
\[
  N'/H
  =\operatorname{Norm}_G(M_\ell)/
   \operatorname{Centr}_G(M_\ell)
\]
by an \(S\)-group \(K'\). Since \(N'\) is smooth and
\(N'_s=N_s\), it is an open subgroup of \(N\) containing \(H\)
(Exposé VI\textsubscript{B}~\SGARefLink{sga3:xv:lemma:2:5}{2.5}). Its image over \(U\) is an open
subgroup of \(K\), isomorphic to \(K'|_U\). Glue \(K\) and \(K'\)
along this open subgroup to obtain \(L\), and glue the projections
\(N|_U\to K\) and \(N'\to K'\) to obtain \(p:N\to L\). Then
\((L,p)\) represents \(N/H\).

\medskip\noindent\textit{Proof of (v).}
First suppose \(S=\operatorname{Spec}k\). The image \(H'\) of \(H\)
is a smooth subgroup of \(G'\). The required assertion follows from:

\SGAStatement{Lemma 7.2}{Let \(u:G\to G'\) be an epimorphism of
smooth connected algebraic groups over a field \(k\), let \(T\) be a
torus of \(G\), and let \(T'\) be its image in \(G'\). Then
\[
 u(\operatorname{Centr}_G(T))
 =\operatorname{Centr}_{G'}(T').
\]}
\SGARefTarget{sga3:xv:lemma:7:2}

Put \(H=\operatorname{Centr}_G(T)\),
\(H'=\operatorname{Centr}_{G'}(T')\), and \(H''=u(H)\). We have
\(H''\subseteq H'\). After extending the base field, assume \(k\)
algebraically closed. It suffices to prove that every Cartan subgroup
of \(H'\) lies in \(H''\): then \(H''\) contains the open subset of
regular points of \(H'\), so is open, and equals the connected group
\(H'\).

Let \(C'\) be a Cartan subgroup of \(H'\). It is also a Cartan
subgroup of \(G'\), since \(H'\), being the centralizer of the torus
\(T'\), has the same reductive and nilpotent ranks as \(G'\). Put
\[
  K=(u^{-1}(C'))^0_{\mathrm{red}}.
\]
Because \(T'\) is central in \(H'\), \(C'\) contains \(T'\), and
\(K\) contains \(T\). Let \(C\) be a Cartan subgroup of \(K\)
containing \(T\). The torus \(T\) lies in the unique maximal torus of
\(C\), central in \(C\) (Exposé \SGARefLink{sga3:xii:6.6:item:c}{XII~6.6(c)}); hence
\(C\subseteq H\). Any two Cartan subgroups of \(G\) are conjugate,
and the image of a Cartan subgroup of \(G\) is a Cartan subgroup of
\(G'\) (Exposé \SGARefLink{sga3:xii:6.6}{XII~6.6}). It follows that \(C\) is a Cartan subgroup
of \(G\), and its image is a Cartan subgroup of \(G'\) contained in
\(C'\). Thus \(u(C)=C'\), proving \(C'\subseteq H''\).

This proves (v) over a field. In general, finite presentation and the
usual reductions permit us to assume \(S\) affine noetherian, then
local, and finally the spectrum of a complete noetherian local ring
\(A\), by fpqc descent of subpreschemes of \(G'\).

Use the notation of (ii). Let \(T_s\) be the maximal central torus of
\(H_s\), and let \(M_\ell\subset H\) lift \({}_\ell T_s\), with
\(H=\operatorname{Centr}_G(M_\ell)\). Let \(T'_s=u_s(T_s)\).
Since \(G'\) is separated, \(u(M_\ell)\) is represented by a subgroup
\(M'_\ell\) of multiplicative type (Exposé \SGARefLink{sga3:ix:theorem:6:8}{IX~6.8}). Put
\[
  H'=\operatorname{Centr}_{G'}(M'_\ell),
\]
which is smooth. For every power \(\ell'\) of \(q\), a sufficiently
large \(\ell\) makes \((M'_\ell)_s\) contain
\({}_{\ell'}T'_s\). We may therefore choose \(\ell\) such that
\[
 H'_s=\operatorname{Centr}_{G'_s}(T'_s)=u_s(H_s),
\]
the last equality being \SGARefLink{sga3:xv:lemma:7:2}{Lemma 7.2}.

The restriction \(v:H\to H'\) of \(u\) is flat. Indeed, \(H,H'\) are
flat over \(S\), \(H_s\to H'_s\) is flat, and hence \(v\) is flat
near \(H_s\) (EGA IV~11.3.10 and 11.3.1). It is therefore flat on an
open subgroup of \(H\) (Exposé VI\textsubscript{B}~2.2), and thus
everywhere because the fibres of \(H\) are connected. The
set-theoretic image of \(H\) is consequently an open subset of \(H'\),
necessarily \((H')^0\), and with its induced structure represents
the fpqc image sheaf \(u(H)\). \SGARefLink{sga3:xv:lemma:7:2}{Lemma 7.2} shows that it belongs to
\(\mathcal{CT}_{G'}(S)\).

If \(H=\operatorname{Centr}_G(T)\), then \(T'=u(T)\) is a sub-torus
of \(G'\) (Exposé \SGARefLink{sga3:ix:theorem:6:8}{IX~6.8}). The image \(u(H)\) is contained in
\(\operatorname{Centr}_{G'}(T')\), agrees with it fibrewise by
\SGARefLink{sga3:xv:lemma:7:2}{Lemma 7.2}, and is smooth. Hence
\[
  u(H)=\operatorname{Centr}_{G'}(T').
\]

\medskip\noindent\textit{Proof of (vi).}
Since \(\mathcal{CT}_G,\mathcal{CT}_{G'}\) are smooth over \(S\), it
is enough to prove faithful flatness on geometric fibres. Thus take
\(S=\operatorname{Spec}k\), with \(k\) algebraically closed. Let
\(H'\in\mathcal{CT}_{G'}(k)\), let \(T'\) be its maximal central
torus, choose a sub-torus \(T\subset G\) mapping onto \(T'\), and put
\[
 H=\operatorname{Centr}_G(T),\qquad
 N=\operatorname{Norm}_G(H),\qquad
 N'=\operatorname{Norm}_{G'}(H').
\]
By \SGARefLink{sga3:xv:lemma:7:2}{Lemma 7.2}, \(H'=u(H)\). Part (iv) identifies \(G/N\) and \(G'/N'\)
with open neighbourhoods of \(H\) and \(H'\) in the two
\(\mathcal{CT}\)-spaces. On these neighbourhoods, \(\widetilde u\)
is
\[
  w:G/N\longrightarrow G'/N',
\]
obtained from \(u\) by passage to quotients. This is an epimorphism
of homogeneous spaces under \(G\), hence faithfully flat. Therefore
\(\widetilde u\) is flat and its image is an open subset of
\(\mathcal{CT}_{G'}\) containing all \(k\)-points. Since
\(\mathcal{CT}_{G'}\) is of finite type, \(\widetilde u\) is
surjective and hence faithfully flat. If \(\ker u\) is central, the
remaining assertions are Exposé \SGARefLink{sga3:xii:7.12}{XII~7.12}.

The next theorem generalizes Exposé \SGARefLink{sga3:xii:7.1}{XII~7.1}.

\SGAStatement{\SGAInlineRefTarget{sga3:xv:theorem:7-3}Theorem 7.3}{Let \(S\) be a prescheme and let \(G\)
be an \(S\)-group prescheme of finite presentation, smooth over \(S\),
with connected fibres. Define the functor
\[
 \mathcal C(S')=\{\text{Cartan subgroups of }G_{S'}\}.
\]
\begin{enumerate}
\renewcommand{\labelenumi}{\textup{\roman{enumi})}}
\item The following are equivalent:
\begin{enumerate}
\renewcommand{\labelenumii}{\textup{\alph{enumii})}}
\item \(\mathcal C\) is representable;
\item \(\mathcal C\) is represented by an \(S\)-prescheme smooth,
quasi-projective, and of finite presentation over \(S\), with affine
fibres;
\item \(G\) admits a Cartan subgroup locally for the étale topology;
\item \(G\) admits a Cartan subgroup locally for the faithfully flat
topology;
\item the nilpotent rank of the fibres of \(G\) is locally constant
on \(S\).
\end{enumerate}
\item If these conditions hold, any two Cartan subgroups of \(G\)
are locally conjugate for the étale topology. The regular locus of
the fibres of \(G\) (Exposé \SGARefLink{sga3:xiii:corollary:2:7}{XIII~2.7}) is an open subprescheme
\(G_{\mathrm{reg}}\), of finite presentation over \(S\), and every
section of \(G_{\mathrm{reg}}\) over \(S\) is contained in a unique
Cartan subgroup of \(G\).
\item Let \(G'\) be an \(S\)-group prescheme of finite presentation
and \(u:G\to G'\) faithfully flat. Then, for every Cartan subgroup
\(C\) of \(G\), its image \(C'=u(C)\) is represented by a Cartan
subgroup of \(G'\), and \(C\to C'\) is faithfully flat.
\item Under the conditions of (i) and (iii), the morphism
\[
 \widetilde u:\mathcal C_G\longrightarrow\mathcal C_{G'},
 \qquad C\longmapsto u(C),
\]
is faithfully flat. If \(\ker u\) is central, it is an isomorphism.
\item Further information on transporters and relations with maximal
tori is given in \SGARefLink{sga3:xv:proposition:7-1}{Proposition 7.1} and Exposé \SGARefLink{sga3:xii:7.1}{XII~7.1}.
\end{enumerate}}
\SGARefTarget{sga3:xv:theorem:7:3}

\SGAProof{Proof}
For (i) we prove
\[
 (\mathrm b)\Rightarrow(\mathrm c)\Rightarrow(\mathrm d)
 \Rightarrow(\mathrm e)\Rightarrow(\mathrm b)
 \Rightarrow(\mathrm a)\Rightarrow(\mathrm d).
\]
Assume (b), and let \(s\in S\). Since \(\mathcal C_s\) is smooth
over \(\kappa(s)\), it has points with residue field a finite
separable extension of \(\kappa(s)\). Exposé \SGARefLink{sga3:xi:proposition:1:10}{XI~1.10} gives an open
neighbourhood \(U\) of \(s\) and an étale surjective morphism
\(U'\to U\) such that \(G_{U'}\) admits a Cartan subgroup. Thus
(b)\(\Rightarrow\)(c), while (c)\(\Rightarrow\)(d) is clear.

Assume (d), and let \(s\in S\). There is a flat \(S\)-prescheme
\(S'\), whose image contains \(s\), such that \(G_{S'}\) admits a
Cartan subgroup. For \(s'\in S'\) above \(s\), the nilpotent rank of
the fibres of \(G_{S'}\) is constant on
\(\operatorname{Spec}\mathcal O_{S',s'}\). Hence the nilpotent rank
of the fibres of \(G\) is constant on
\(\operatorname{Spec}\mathcal O_{S,s}\), which is the image of that
local spectrum (EGA IV~2.3.4(ii)). Let the value be \(r\). By 6.3 bis,
the set \(E_r\) where the nilpotent rank equals \(r\) is
ind-constructible; it contains a neighbourhood of \(s\) by
EGA IV~1.10.1. Thus (d)\(\Rightarrow\)(e).

Assume (e). Locally on \(S\), the nilpotent rank is a constant \(r\).
For every \(S'\), the Cartan subgroups of \(G_{S'}\) are exactly the
smooth subgroup preschemes of relative dimension \(r\) whose
geometric fibres are centralizers of tori. By \SGARefLink{sga3:xv:proposition:7-1}{Proposition 7.1},
\(\mathcal C\) is the open-and-closed subprescheme of
\(\mathcal{CT}\) representing the relative-dimension-\(r\)
subfunctor. It has all the properties in (b); affineness of its
fibres is Exposé \SGARefLink{sga3:xii:theorem:7-1:item:d}{XII~7.1(d)}. This proves
(e)\(\Rightarrow\)(b), and (b)\(\Rightarrow\)(a) is immediate.

Finally assume (a). By \SGARefLink{sga3:xv:proposition:6-2:item:iii}{6.2(iii)}, \(\mathcal C\) commutes with filtered
inductive limits of rings, so its representing prescheme is locally
of finite presentation (EGA IV~8.14.2). To prove it smooth, it is
enough to show that, when \(S\) is affine and \(S_0\subseteq S\) is
defined by a nilpotent ideal, every Cartan subgroup \(C_0\) of
\(G_0\) lifts to a Cartan subgroup of \(G\). Existence of \(C_0\)
implies (e) on \(S_0\), hence on \(S\), which has the same underlying
space. Since (e)\(\Rightarrow\)(b), \(\mathcal C\) is smooth.
It is also surjective over \(S\), so (a)\(\Rightarrow\)(d).

For (ii), if \(C,C'\) are Cartan subgroups, then
\(\operatorname{Transp}_G(C,C')\) is represented by a smooth
\(S\)-prescheme (\SGARefLink{sga3:xv:proposition:7-1:item:iii}{Proposition 7.1(iii)}) with nonempty fibres
(Exposé \SGARefLink{sga3:xii:theorem:6-6:item:a}{XII~6.6(a)},(c)). Hensel's lemma (Exposé \SGARefLink{sga3:xi:proposition:1:10}{XI~1.10}) gives local
conjugacy for the étale topology. The assertions about regular points
follow from Exposé \SGARefLink{sga3:xiii:theorem:3:1}{XIII~3.1} and \SGARefLink{sga3:xiii:corollary:3:2}{3.2}.

For (iii), \SGARefLink{sga3:xv:proposition:7-1:item:v}{Proposition 7.1(v)} represents \(u(C)\) by a smooth
subgroup \(C'\) of \(G'\). Its fibres are Cartan subgroups of the
fibres of \(G'\) (Exposé \SGARefLink{sga3:xii:6.6:item:d}{XII~6.6(d)}), so \(C'\) is a Cartan subgroup.

For (iv), faithful flatness is proved as in \SGARefLink{sga3:xv:proposition:7-1:item:vi}{Proposition 7.1(vi)}.
If \(\ker u\) is central and \(C'\) is a Cartan subgroup of \(G'\),
then \(C=u^{-1}(C')\) is smooth with connected fibres
(\SGARefLink{sga3:xv:proposition:7-1:item:vi}{Proposition 7.1(vi)}), and its fibres are Cartan subgroups (Exposé
\SGARefLink{sga3:xii:6.6:item:e}{XII~6.6(f)}). Thus \(C\) is a Cartan subgroup, and \SGARefLink{sga3:xv:proposition:7-1:item:vi}{Proposition
7.1(vi)} completes the proof.
